\section{A canonical detected quotient of class five}
\label{sec:canonical-class-five}

Retain $J_5$, $K_4^{\rm det}$ and the class-four quotient.
Put
\[
 J_6=\overline{[H,K_4^{\rm det}][B,K_3][K_2,K_2]},
 \qquad R_5=K_4^{\rm det}/J_6,
 \qquad S_5=J_5/J_6,
\]
let \(K_5^{\rm det}\) be the inverse image of \((R_5)_{\Lambda\text{-tor}}\) in \(K_4^{\rm det}\), and put \(Q_5=H/K_5^{\rm det}\). For \(2\leq j\leq5\), write \(L^{(j)}\) for the fixed field of the kernel of \(H\to Q_j\). We abbreviate
\[
 E_2=B/K_2=\bar E,
 \qquad E_3=K_2/K_3=R_3,
 \qquad E_4=K_3/K_4^{\rm det}.
\]
These have parities even, odd, even, respectively.
The three-subgroup identity proves $J_6\subseteq J_5$
and centrality of $R_5,S_5$ in $H/J_6$. The exact
quotient is
\begin{equation}\label{eq:class5-old-torsion}
 0\longrightarrow S_5\longrightarrow R_5
                  \longrightarrow(R_4)_{\rm tor}\longrightarrow0.
\end{equation}
It retains all of the preceding bounded $p$-torsion.



\begin{theorem}[Canonical detection through class five]\label{thm:canonical-class-five}
For every odd prime $p$, the $\Lambda$-torsion submodule
$(R_5)_{\rm tor}$ is closed. Its inverse image
$K_5^{\rm det}$ in $K_4^{\rm det}$ is fully
Galois-stable and contains the outer kernel.
The quotient $Q_5=H/K_5^{\rm det}$ is torsionfree of
nilpotency class exactly five and is the compact
image in the full rational outer target through
generator length five. It fits into the central
extension
\[
 1\longrightarrow R_5/(R_5)_{\rm tor}
 \longrightarrow Q_5\longrightarrow Q_4\longrightarrow1.
\]
The natural map
$\operatorname{Lie}^5(W)\to R_5/(R_5)_{\rm tor}$ is injective.
Together with its preceding central quotients it gives
\[
 L^{(2)}\subsetneq L^{(3)}\subsetneq L^{(4)}
                         \subsetneq L^{(5)}\subseteq\Omega_p^*.
\]
\end{theorem}

\begin{proof}[Proof of Theorem~\ref{thm:canonical-class-five}]
The central module is computed using a commutator
surjection that becomes an isomorphism after inverting $p$.
The commutators induce the surjection
\[
 A_5=(W\widehat\otimes E_4)\oplus
                              (E_2\widehat\otimes E_3)
                         \longrightarrow S_5.
\]
The three factors of $J_6$ give precisely the
independence of lifts in its two summands.
Centrality gives bilinearity. The Hall--Witt relations
of types $(1,1,3)$ and $(1,2,2)$ are
\begin{align*}
 j_{113}((x\wedge y)\otimes r)
   &=\bigl(x\otimes[y,r]-y\otimes[x,r],
                                  -\nu(x\wedge y)\otimes r\bigr),\\
 j_{122}(x\otimes(e\wedge f))
   &=\bigl(x\otimes[e,f],
                                  f\otimes[x,e]-e\otimes[x,f]\bigr).
\end{align*}
Here all lower bracket maps are those already defined
in $Q_4$. With the closed relation image, put
\begin{equation}\label{eq:class5-candidate}
 C_5=A_5/\operatorname{im}(j_{113},j_{122}).
\end{equation}
We have a canonical surjection $C_5\twoheadrightarrow S_5$,
and both modules are odd under $c$.

We first prove that this surjection is an isomorphism
after inverting $p$. There is a concrete reason to
make this passage explicitly: the odd Hopf quotient
at this stage can include the earlier bounded module
$S_3[p^\infty]$. It cannot be discarded integrally.
Since $\gamma_5H\subseteq K_4^{\rm det}$, we have
$\gamma_6H\subseteq J_6$. Thus the calculation takes
place in fixed ordinary nilpotency class five.
The rational BCH and inverse-substitution formulas
have only finitely many fixed denominators there.
To justify the bounded completion, use the Magnus
embedding in the completed tensor algebra on
$H^{\rm ab}$. The formulas for $\log(1+X)$, $\exp(X)$
and multiplication through ordinary degree five have
the same finite list of rational coefficients at
every finite rank. One integer clears these
coefficients independently of the number of basis
vectors. Passing to the inverse limit over finite
coordinate sets therefore gives the free rational
nilpotent Lie algebra on the compact module $M$,
with a common bound on every coefficient denominator.

The abelian images of $B,K_2,K_3,K_4^{\rm det}$
are, respectively,
\[
 D\oplus T,\quad D\oplus\pi Z,\quad
 D_0\oplus\pi Z,\quad D_0\oplus Z_0.
\]
Choose rational complements in this finite-dimensional
flag and compatible representatives, assigning weights
\[
 W_L:1,\quad (T/\pi Z)_L:2,\quad (D/D_0)_L:3,
 \quad(\pi Z/Z_0)_L:4,\quad(D_0\oplus Z_0)_L:5.
\]
For $W$ keep a compact $\mathbf Z_p$-basis and lifts.
The remaining choices involve finitely many vectors
and can be made after a fixed bounded scaling.
Substitution by these representatives in the rational
free nilpotent Lie algebra of $H$ is invertible:
its linear term is invertible, and its inverse is
constructed successively through the five ordinary
lower-central steps.

The proved presentations of $Q_2,Q_3,Q_4$ identify
their rational pieces with the free weighted pieces
through weights two, three, four. The map from the
free weighted truncation to $Q_4$, using the compatible
representatives, is therefore an isomorphism on every
graded piece and hence on the finite filtered object.
It identifies the maps to $Q_3,Q_2,W$ with the
corresponding truncations. In the rational free source,
the images of $H,B,K_2,K_3,K_4^{\rm det}$ are exactly
the ideals of weight at least $1,2,3,4,5$.

The image of $J_6$ is consequently all commutators of
weight at least six: every such bracket lies in one
of the three displayed factors. The image of $J_5$
is all commutators of weight at least five. Their
quotient is exactly
\[
 L_5(W_L,T',D',Z'),\qquad
 T'=(T/\pi Z)_L,\quad D'=(D/D_0)_L,\quad
 Z'=(\pi Z/Z_0)_L,
\]
with generator weights $1,2,3,4$.
The weight-five generators themselves do not occur
in this derived quotient. The rational candidate
\eqref{eq:class5-candidate} is the same free weighted
piece: its central bracket types are $(1,4),(2,3)$,
and its new Jacobi types are exactly $(1,1,3),(1,2,2)$.
All lower relations have already been imposed.
The commutator map therefore gives
\begin{equation}\label{eq:class5-rational-source}
              C_5[1/p]\xrightarrow{\ \sim\ }S_5[1/p].
\end{equation}
In particular its integral kernel is $p$-torsion.

We next bound that torsion uniformly. Retain $b,c$
from the class-four proof and put
$d_4=\max\{b+c,2b\}$. Choose $c_4$ killing the finite
torsion subgroup of $\pi Z/Z_0$. The earlier weighted
Hall lattices satisfy
\[
 E_{2,0}=N\oplus C_2\subset E_2,\quad p^bE_2\subset E_{2,0},
\]
\[
 E_{3,0}=L_3(W,C_2)\oplus C_3\subset E_3,
                      \quad p^{b+c}E_3\subset E_{3,0}.
\]
The exact sequence
$0\to S_4/S_4[p^\infty]\to E_4\to\pi Z/Z_0\to0$
and the class-four Hall comparison give a finite
free section $C_4$ and
\[
 E_{4,0}=L_4(W,C_2,C_3)\oplus C_4\subset E_4,
                 \qquad p^{d_4+c_4}E_4\subset E_{4,0}.
\]
These are compact $\mathbf Z_p$-lattices; no integral
$\Lambda$-freeness or $\Lambda$-stability is required.

Set
\[
 A_{5,0}=(W\widehat\otimes E_{4,0})\oplus
                          (E_{2,0}\widehat\otimes E_{3,0}),
 \qquad d_5=\max\{d_4+c_4,2b+c\}.
\]
Then $p^{d_5}A_5\subseteq A_{5,0}$. Let $I_0$ be the
two Jacobi images on these free weighted lattices.
The quotient $A_{5,0}/I_0=L_5(W,C_2,C_3,C_4)$ is a
free Hall module and is $p$-torsionfree. If $I$ is
the full relation image in \eqref{eq:class5-candidate},
then
\[
                              p^{d_4}I\subseteq I_0:
\]
the $(1,1,3)$ relation needs the factor $p^{b+c}$,
and the $(1,2,2)$ relation needs $p^{2b}$.
For $p^ka\in I$, the element $p^{d_5}a$ lies in
$A_{5,0}$ and has a $p$-power multiple in $I_0$.
Saturation of $I_0$ there gives $p^{d_5}a\in I_0$.
Thus $p^{d_5}C_5[p^\infty]=0$, and
$A_{5,0}/I_0\hookrightarrow C_5$ as well: its kernel
is killed by $p^{d_4}$ and is therefore zero.
By \eqref{eq:class5-rational-source}, every $p$-torsion
element of $S_5$ lifts to $p$-torsion in $C_5$.
We have proved
\begin{equation}\label{eq:class5-pbound}
 p^{d_5}S_5[p^\infty]=0,
 \qquad\operatorname{Lie}^5(W)\hookrightarrow S_5.
\end{equation}
This argument keeps the possible integral kernel of
$C_5\to S_5$ instead of asserting that it is zero.

For rational detection, the new finite quotient is
$Z'$ of weight four. Its companions lie in
$S_{4,L}=L_4(W_L,T',D')$. We verify the total-product
divisibility needed for their residues. Let $A^0$
be the pair algebra of
Theorem~\ref{thm:pair-total-saturation}, and let $A$
also contain the cubic generators $D'$. In degree
four there is an exact $\Gamma$-stable sequence
\[
 0\longrightarrow A^0_4\longrightarrow A_4
 \longrightarrow(W_L\widehat\otimes D')
                    \oplus(D'\widehat\otimes W_L)\longrightarrow0.
\]
The quotient action is diagonal, and it has no
$(\gamma-\chi)$-invariants by the regular singleton
factor and its finite Jordan determinant.
On $z_1z_2z_3z_4=\chi$, a prefix triple pole at $\beta$
becomes $z_4-\chi/\beta$; a suffix triple pole becomes
$z_1-\chi/\beta$. These singleton divisors are distinct.
Every pair denominator, including an internal
companion denominator, is a unit at each generic
singleton divisor, as is the other triple divisor.

The Jordan cokernel in a singleton variable $u$ is
that of $uA_D-\chi$, namely the coefficient module
modulo $(u\beta-\chi)^r$ for a block of length $r$.
Its kernel image has the successive principal parts
$C_3/(\chi/u-\beta)^{r-j}$. The already checked total
degree-three residue independence, with a finite
color and coefficient minor in the internal triple
variables, detects every singleton jet. A zero total
image therefore makes the quotient class a
$(\gamma-\chi)$-multiple. Subtract it in the source.
The remainder belongs to $A^0_4$, where
Theorem~\ref{thm:pair-total-saturation} supplies the
division. Adding proper triple denominators does
not change total-divisor membership for this
remainder: each is coprime to the prime total divisor.
Iteration gives every power, and the Eulerian
projector gives the same assertion on $S_{4,L}$.

The $\Lambda_L$-torsionfreeness of $E_{4,L}$ now makes
the $Z'$ Jordan companions independent in
$S_{4,L}/(\gamma-\alpha)S_{4,L}$; the preceding
divisibility transports this to target
residues. Their kernels have a full four-position
denominator and only proper pair or triple poles
in the numerator. To detect the free weighted
degree-five source, extract an initial four-position
pole first, then an initial triple pole, then an
initial pair pole, and finally a singleton. Every
residue coefficient involves only disjoint tail
variables; the finite constant extraction matrices
and induction give an injective associative map.
It restricts to the free Lie piece in
\eqref{eq:class5-rational-source}. This proves
\begin{equation}\label{eq:class5-rational-injection}
 S_5[1/p]\hookrightarrow\text{target generator length five}.
\end{equation}
In particular $S_5[1/p]$ is $\Lambda$-torsionfree,
because a nonzero $f(\gamma)$ acts on its coefficient
arrays by the nonzero multiplier $f(z_1\cdots z_5)$.

The preceding exact formula kills $(R_4)_{\rm tor}$
by $p^{b+d_4}a(\gamma)$. Equations
\eqref{eq:class5-old-torsion},
\eqref{eq:class5-pbound} and
\eqref{eq:class5-rational-injection} give
\begin{equation}\label{eq:class5-closed-torsion}
 (R_5)_{\rm tor}
      =\ker\bigl(p^{b+d_4+d_5}a(\gamma):R_5\to R_5\bigr).
\end{equation}
Indeed a torsion element multiplied by
$p^{b+d_4}a(\gamma)$ lies in the rationally torsionfree
module $S_5$ and hence is killed by $p^{d_5}$.
The converse follows from the nonzero annihilator.
This proves closedness and full arithmetic invariance.

The exact class-four target kernel is $K_4^{\rm det}$.
The three factors of $J_6$ therefore have target
length at least six, so the image of $R_5$ is in pure
length five. For an outer-kernel element,
multiply its class in $R_5$ by $p^{b+d_4}a(\gamma)$.
The result lies in $S_5$ and has zero rational image;
\eqref{eq:class5-pbound} then proves that its class
lies in \eqref{eq:class5-closed-torsion}.
Conversely the roots of $a$ are among the original
abelian torsion eigenvalues. By
Lemma~\ref{lem:positive-eigenvalues}, none occurs on
a positive target coordinate of length five.
Thus the displayed torsion subgroup itself maps to
zero. This proves exactness of the compact-image map.

The central extension is torsionfree since $Q_4$ is
torsionfree and its central kernel has no $p$-torsion.
The pure Hall module in \eqref{eq:class5-pbound}
has an injective five-variable coefficient map,
so meets $(R_5)_{\rm tor}$ trivially. The word
$[x,[x,[x,[x,y]]]]$ is nonzero for independent odd
coordinates. Finally $\gamma_6H\subseteq J_6$,
which proves class exactly five and strictness.
\end{proof}

Put
\[
 D_1=\operatorname{im}\bigl((R_5)_{\rm tor}^-
                                         \longrightarrow D_0\bigr).
\]
Since $S_5$ is odd, the even part of $R_5$ is entirely
$\Lambda$-torsion and has image $Z_0$ in $M$.
The abelianization map therefore gives
\begin{equation}\label{eq:class5-abelianization}
                         Q_5^{\rm ab}=(U/D_1)\oplus(T/Z_0).
\end{equation}

\subsection{Selected components at class five}
\label{subsec:selected-class-five}

Let $B_5$ be the image of $B$ in $Q_5$. It still has
class at most two, since
$[[B,B],B]\subseteq[B,K_3]\subseteq J_6$.
Thus its integral class-two BCH Lie algebra exists
for every odd $p$. For a good even set $\Sigma$, take
the closed $Q_5$-normal subgroup generated by the
exponentials of all omitted even isotypic submodules
of $\operatorname{Lie}(B_5)$, and denote the quotient
by $Q_{5,\Sigma}$. Write $J_\Sigma^{(5)}$ for this closed normal subgroup, and $L_\Sigma^{(5)}$ for the fixed field of $\ker(H\to Q_{5,\Sigma})$.



\begin{corollary}[Class-five component detection]
\label{cor:selected-class-five}
Let $p$ be odd and let $\Sigma$ be any set of even exponents
such that $X^{[e]}$ is finite for every $e\in\Sigma$.
The subgroup $J_\Sigma^{(5)}$ is independent of the tame
complement and stable under the full arithmetic Galois group.
Its canonical field is Galois over $\mathbb Q$, and
\[
 L_\Sigma^{(5)}\subseteq\Omega_p^*,\qquad
 \operatorname{Gal}(L_\Sigma^{(5)}/F_\infty)=Q_{5,\Sigma},\qquad
 Q_{5,\Sigma}^{\rm ab}=(U/D_1)\oplus T_\Sigma.
\]
For $-2,-4\in\Sigma$, the quotient has nilpotency class
exactly five. It maps onto the class-four component
quotient of Proposition~\ref{prop:selected-class-four},
and strictly enlarges its field in this case.
\end{corollary}

\begin{proof}[Proof of Corollary~\ref{cor:selected-class-five}]
Tame complements are $H$-conjugate, and for each
arithmetic element $g$ a suitable $h^{-1}g$ centralizes
the complement pointwise. The same normal-closure
argument as at class four proves independence and
full arithmetic stability. The image of the omitted
even Lie submodules in $B_4$ is exactly the earlier
set of generators, giving the stated surjection.
Their abelian image in \eqref{eq:class5-abelianization}
is precisely the omitted even part of $T/Z_0$.
For selected $e$, the component theorem gives
$(Z_0)_e=0$, proving the displayed abelianization.

Project the free target through length five
onto two distinct odd generators of weights
$m,n\equiv-1\pmod{p-1}$, after the same
$\Delta$-equivariant choice of lifts as before.
The image of $B_5$ has no linear term, so its even
part has only character $-2$ in length two and
character $-4$ in length four. The morphism therefore
kills every omitted even generator and factors
through $Q_{5,\Sigma}$. The two moment coordinates
of $w,\gamma w$ have the nonzero determinant
$\lambda_m\lambda_n(u^n-u^m)$; their fifth Hall
commutator has nonzero length-five term. It survives
in the selected quotient, proving class five and
strictness.
\end{proof}
